Về phạm vi nghiên cứu và ứng dụng: Đề tài tập trung ứng dụng mô hình Copula như một công cụ cốt lõi trong quản trị rủi ro danh mục đầu tư. Phương pháp này cho phép tách biệt phân phối biên của từng cổ phiếu khỏi cấu trúc phụ thuộc chung (mô hình hóa từ dưới lên). Nhờ tính linh hoạt trên thang đo phân vị và dễ dàng kết hợp với mô phỏng Monte Carlo, nghiên cứu có thể giải quyết các bài toán định lượng rủi ro đuôi (Value at Risk - Var) và dự báo các sự kiện cực đoan đa biến khi thị trường chứng khoán xảy ra khủng hoảng đồng loạt.\\
Về giới hạn của nghiên cứu: Việc ứng dụng mô hình này cũng đối mặt với một số thách thức nhất định. Rủi ro lớn nhất đến từ việc lựa chọn sai cấu trúc Copula (rủi ro mô hình), có thể dẫn đến đánh giá thấp các thảm họa mang tính hệ thống. Bên cạnh đó, mô hình Copula truyền thống được sử dụng trong bài thường mang tính tĩnh, thiếu linh hoạt khi mở rộng trong không gian đa chiều, và đặc biệt khó hiệu chuẩn chính xác do sự khan hiếm của các điểm dữ liệu lịch sử về những đợt tổn thất cực đoan đồng thời trên thị trường.